\documentclass[12pt]{article}
\usepackage[a4paper,margin=1in]{geometry}
\usepackage[utf8]{inputenc}
\usepackage[T1]{fontenc}
\usepackage{amsmath,amssymb}
\usepackage{textcomp}
\usepackage{parskip}
\begin{document}
\begin{center}
{\Large \textbf{2001 STEP III}}
\end{center}

\subsection*{Page 1}
\begin{flushleft}
STEP III, 2001 2\par
\bigskip\textbf{Section A: Pure Mathematics}\par
dy 1\par
\bigskip\textbf{Question 1.} Given that y = ln x + \sqrt(x2 + 1) , show that = .\par
\[dx \sqrt(x^2 + 1)\]
\[Prove by induction that, for n > 0 ,\]
\[x^2 + 1 y(n+2) + (2n + 1) xy(n+1) + n2y(n) = 0 ,\]
n\par
d y\par
\[where y(n) = and y(0) = y .\]
dxn\par
Using this result in the case x = 0 , or otherwise, show that the Maclaurin series for y begins\par
\[x^3 3x5\]
x +\par
- 6 40\par
and find the next non-zero term.\par
\bigskip\textbf{Question 2.} Show that cosh -1 x = ln (x + \sqrt(x2 1)) .\par
-\par
\[Show that the area of the region defined by the inequalities y^2 > x^2 8 and x^2 > 25y2 16\]
- -\par
\[is (72/5) ln 2.\]
\bigskip\textbf{Question 3.} Consider the equation\par
\[x^2 bx + c = 0 ,\]
-\par
where b and c are real numbers.\par
\textbf{(i)} Show that the roots of the equation are real and positive if and only if b > 0 and\par
\[b^2 > 4c > 0, and sketch the region of the b -c plane in which these conditions hold.\]
\textbf{(ii)} Sketch the region of the b -c plane in which the roots of the equation are real and less\par
than 1 in magnitude.\par
\end{flushleft}
\subsection*{Page 2}
\begin{flushleft}
STEP III, 2001 3\par
\bigskip\textbf{Question 4.} In this question, the function sin -1 is defined to have domain 1 6 x 6 1 and range\par
\[\pi 6 x 6 \pi and the function tan -1 is defined to have the rea - l numbers as its domain\]
\[- 2 \pi 2 \pi\]
\[and range < x < .\]
- 2 2\par
\textbf{(i)} Let\par
2x\par
\[g (x) = , < x < .\]
\[1 + x^2 -\infty \infty\]
Sketch the graph of g (x) and state the range of g .\par
\textbf{(ii)} Let\par
2x\par
\[f (x) = sin\]
-1\par
\[, < x < .\]
\[\lambda 1 + x^2 \mu -\infty \infty\]
\[Show that f (x) = 2 tan -1 x for 1 6 x 6 1 and f (x) = \pi 2 tan -1 x for x > 1 .\]
- -\par
Sketch the graph of f (x).\par
\bigskip\textbf{Question 5.} Show that the equation x3 + px + q = 0 has exactly one real solution if p > 0 .\par
A parabola C is given parametrically by\par
\[x = at2, y = 2at (a > 0) .\]
Find an equation which must be satisfied by t at points on C at which the normal passes\par
through the point (h, k) . Hence show that, if h 6 2a , exactly one normal to C will pass\par
through (h, k) .\par
Find, in Cartesian form, the equation of the locus of the points from which exactly two normals\par
can be drawn to C . Sketch the locus.\par
\bigskip\textbf{Question 6.} The plane\par
x y z\par
\[+ + = 1\]
a b c\par
meets the co-ordinate axes at the points A, B and C . The point M has coordinates 1 a, 1 b, 1 c\par
2 2 2\par
and O is the origin.\par
Show that OM meets the plane at the centroid 1 a, 1 b, 1 c of triangle ABC. Show also that\par
3 3 3\par
the perpendiculars to the plane from O and from M meet the plane at the orthocentre and\par
at the circumcentre of triangle ABC respectively.\par
Hence prove that the centroid of a triangle lies on the line segment joining its orthocentre and\par
circumcentre, and that it divides this line segment in the ratio 2 : 1 .\par
[The orthocentre of a triangle is the point at which the three altitudes intersect; the circum-\par
centre of a triangle is the point equidistant from the three vertices.]\par
\end{flushleft}
\subsection*{Page 3}
\begin{flushleft}
STEP III, 2001 4\par
\bigskip\textbf{Question 7.} Sketch the graph of the function ln x 1 x2.\par
- 2\par
Show that the differential equation\par
dy 2xy\par
\[=\]
\[dx x^2 1\]
-\par
describes a family of parabolas each of which passes through the points (1 , 0) and ( 1 , 0)\par
-\par
and has its vertex on the y-axis.\par
Hence find the equation of the curve that passes through the point (1 , 1) and intersects each\par
of the above parabolas orthogonally. Sketch this curve.\par
[Two curves intersect orthogonally if their tangents at the point of intersection are perpendic-\par
ular.]\par
\bigskip\textbf{Question 8.} (i) Prove that the equations\par
\[z (1 + i) 2 = 2 ( )\]
| - | ∗\par
and\par
\[z (1 i) 2 = 2 z 1 2\]
| - - | | - |\par
describe the same locus in the complex z-plane. Sketch this locus.\par
\textbf{(ii)} Prove that the equation\par
\[z 2 \pi\]
\[arg - = ( )\]
\lambda z \mu 4 ∗∗\par
describes part of this same locus, and show on your sketch which part.\par
\textbf{(iii)} The complex number w is related to z by\par
\[w = .\]
z\par
Determine the locus produced in the complex w-plane if z satisfies ( ). Sketch this locus\par
∗\par
and indicate the part of this locus that corresponds to ( ).\par
∗∗\par
\end{flushleft}
\subsection*{Page 4}
\begin{flushleft}
STEP III, 2001 5\par
\bigskip\textbf{Section B: Mechanics}\par
9 B and B are parallel, thin, horizontal fixed beams. B is a vertical distance d sin α above\par
1 2 1\par
\[B , and a horizontal distance d cos α from B , where 0 < α < \pi/2 . A long heavy plank is\]
2 2\par
held so that it rests on the two beams, perpendicular to each, with its centre of gravity at\par
B . The coefficients of friction between the plank and B and B are µ and µ , respectively,\par
1 1 2 1 2\par
\[where µ < µ and µ + µ = 2 tan α .\]
1 2 1 2\par
The plank is released and slips over the beams experiencing a force of resistance from each\par
beam equal to the limiting frictional force (i.e. the product of the appropriate coefficient of\par
friction and the normal reaction). Show that it will come to rest with its centre of gravity\par
over B in a time\par
d 2\par
\[\pi .\]
\lambda g(µ µ ) cos α \mu\par
2 1\par
-\par
\bigskip\textbf{Question 10.} Three ships A, B and C move with velocities v , v and u respectively. The velocities of A\par
1 2\par
and B relative to C are equal in magnitude and perpendicular. Write down conditions that\par
u, v and v must satisfy and show that\par
1 2\par
\[u 1 (v + v ) 2 = 1 (v v ) 2\]
- 2 1 2 2 1 - 2\par
\{ \{ \{ \{\par
\{ \{ \{ \{\par
and\par
\[u 1 (v + v ) (v v ) = 0 .\]
- 2 1 2 · 1 - 2\par
Explain why these equations determine, for given v and v , two possible velocities for C ,\par
1 2\par
\[provided v = v .\]
1 2\par
\[If v and v are equal in magnitude and perpendicular, show that if u = 0 then u = v + v .\]
1 2 1 2\par
\bigskip\textbf{Question 11.} A uniform cylinder of radius a rotates freely about its axis, which is fixed and horizontal.\par
The moment of inertia of the cylinder about its axis is I . A light string is wrapped around\par
the cylinder and supports a mass m which hangs freely. A particle of mass M is fixed to the\par
surface of the cylinder. The system is held at rest with the particle vertically below the axis\par
of the cylinder, and then released. Find, in terms of I, a, M , m, g and θ, the angular velocity\par
of the cylinder when it has rotated through angle θ .\par
Show that the cylinder will rotate without coming to a halt if m/M > sin α , where α satisifes\par
\[α = tan 1 α and 0 < α < \pi .\]
\end{flushleft}
\subsection*{Page 5}
\begin{flushleft}
STEP III, 2001 6\par
\bigskip\textbf{Section C: Statistics}\par
\bigskip\textbf{Question 12.} A bag contains b black balls and w white balls. Balls are drawn at random from the bag and\par
when a white ball is drawn it is put aside.\par
\textbf{(i)} If the black balls drawn are also put aside, find an expression for the expected number\par
of black balls that have been drawn when the last white ball is removed.\par
\textbf{(ii)} If instead the black balls drawn are put back into the bag, prove that the expected\par
number of times a black ball has been drawn when the first white ball is removed is\par
b/w . Hence write down, in the form of a sum, an expression for the expected number\par
of times a black ball has been drawn when the last white ball is removed.\par
\bigskip\textbf{Question 13.} In a game for two players, a fair coin is tossed repeatedly. Each player is assigned a sequence\par
of heads and tails and the player whose sequence appears first wins. Four players, A, B, C and\par
D take turns to play the game. Each time they play, A is assigned the sequence TTH (i.e. Tail\par
then Tail then Head), B is assigned THH, C is assigned HHT and D is assigned HTT.\par
\textbf{(i)} A and B play the game. Let p , p , p and p be the probabilities of A winning\par
HH HT TH TT\par
the game given that the first two tosses of the coin show HH, HT, TH and TT,\par
\[respectively. Explain why p = 1 , and why p = 1 p + 1 p . Show that\]
TT HT 2 TH 2 TT\par
\[p = p = 2 and that p = 1 . Deduce that the probability that A wins the game\]
HH HT 3 TH 3\par
is 2 .\par
\textbf{(ii)} B and C play the game. Find the probability that B wins.\par
\textbf{(iii)} Show that if C plays D, then C is more likely to win than D, but that if D plays A,\par
then D is more likely to win than A.\par
\bigskip\textbf{Question 14.} A random variable X is distributed uniformly on [ 0 , a ]. Show that the variance of X is 1 a2.\par
A sample, X and X , of two independent values of the random variable is drawn, and the\par
1 2\par
\[variance V of the sample is determined. Show that V = 1 (X X )2, and hence prove that\]
4 1 - 2\par
2V is an unbiased estimator of the variance of X.\par
Find an exact expression for the probability that the value of V is less than 1 a2 and estimate\par
the value of this probability correct to one significant figure.\par
\end{flushleft}
\end{document}
